\section{宣统转型其余账本事件锚点}
本节补齐宣统转型余下事件，使咨议、铁路、军政、地方独立、谈判和退位的材料节点全部可导航。
\subsection{1910—1919}
\eventanchor{EQ-1911-WUCHANG-UPRISING-AFTERMATH}{「武昌起义爆发，新军与革命团…}宣统三年（1911年），「武昌起义爆发，新军与革命团体控制武汉部分地区」引发的善后、执行或地方回应构成可由不同材料追踪的后续节点。核心取证：《宣统政纪》宣统三年相关条；军机处档、战事方略及地方奏报；相关地区的档案、志书、契约、报刊或对方材料。该项锚定的是后续追踪问题，影响范围须按地区与群体分别判断。来源保留：战后、改革后或条约后的记忆、地方记录和官方总结常具有事后选择性；后续影响须按地点、群体和时间分层，不作整体化推断。
\eventanchor{EQ-1911-PROVINCIAL-INDEPENDENCE-PRELUDE}{「多省宣布脱离清廷，帝国行政…}宣统三年（1911年），「多省宣布脱离清廷，帝国行政和财政网络快速分裂」之前的条件、信息和争议已通过中央与地方材料显现，构成该事件可追踪的前史节点。核心取证：《宣统政纪》宣统三年相关条；宫中朱批奏折、内阁大库题本与《清史稿》相关志传；同时期地方档案、地方志、日记或对方材料应按具体地区补检。该项锚定的是前史条件与信息背景，不把条件直接写成唯一因果。来源保留：官修编年和地方材料的记录密度与立场并不相同；本行不将条件、传闻、呈报或后见解释直接等同为确定因果。
\eventanchor{EQ-1911-PROVINCIAL-INDEPENDENCE-DECISION}{围绕「多省宣布脱离清廷，帝国…}宣统三年（1911年），围绕「多省宣布脱离清廷，帝国行政和财政网络快速分裂」的命令、调度、照会或呈报形成独立的中央—地方行动文书链。核心取证：《宣统政纪》宣统三年相关条；宫中朱批奏折、内阁大库题本与《清史稿》相关志传。该项锚定的是命令或决策环节，命令抵达和结果仍须另外核验。来源保留：奏折、上谕、部议、军机档或条约可证明行政动作和机构立场，不自动证明所有地区、群体或对手按其意图行动。
\eventanchor{EQ-1911-PROVINCIAL-INDEPENDENCE-AFTERMATH}{「多省宣布脱离清廷，帝国行政…}宣统三年（1911年），「多省宣布脱离清廷，帝国行政和财政网络快速分裂」引发的善后、执行或地方回应构成可由不同材料追踪的后续节点。核心取证：《宣统政纪》宣统三年相关条；宫中朱批奏折、内阁大库题本与《清史稿》相关志传；相关地区的档案、志书、契约、报刊或对方材料。该项锚定的是后续追踪问题，影响范围须按地区与群体分别判断。来源保留：战后、改革后或条约后的记忆、地方记录和官方总结常具有事后选择性；后续影响须按地点、群体和时间分层，不作整体化推断。
\eventanchor{EQ-1912-REPUBLIC-FOUNDING-PRELUDE}{「中华民国临时政府在南京成立…}宣统四年（1912年），「中华民国临时政府在南京成立，与清廷退位谈判并行」之前的条件、信息和争议已通过中央与地方材料显现，构成该事件可追踪的前史节点。核心取证：《宣统政纪》宣统四年相关条；宫中朱批奏折、内阁大库题本与《清史稿》相关志传；同时期地方档案、地方志、日记或对方材料应按具体地区补检。该项锚定的是前史条件与信息背景，不把条件直接写成唯一因果。来源保留：官修编年和地方材料的记录密度与立场并不相同；本行不将条件、传闻、呈报或后见解释直接等同为确定因果。
\eventanchor{EQ-1912-REPUBLIC-FOUNDING-DECISION}{围绕「中华民国临时政府在南京…}宣统四年（1912年），围绕「中华民国临时政府在南京成立，与清廷退位谈判并行」的命令、调度、照会或呈报形成独立的中央—地方行动文书链。核心取证：《宣统政纪》宣统四年相关条；宫中朱批奏折、内阁大库题本与《清史稿》相关志传。该项锚定的是命令或决策环节，命令抵达和结果仍须另外核验。来源保留：奏折、上谕、部议、军机档或条约可证明行政动作和机构立场，不自动证明所有地区、群体或对手按其意图行动。
\eventanchor{EQ-1912-REPUBLIC-FOUNDING-AFTERMATH}{「中华民国临时政府在南京成立…}宣统四年（1912年），「中华民国临时政府在南京成立，与清廷退位谈判并行」引发的善后、执行或地方回应构成可由不同材料追踪的后续节点。核心取证：《宣统政纪》宣统四年相关条；宫中朱批奏折、内阁大库题本与《清史稿》相关志传；相关地区的档案、志书、契约、报刊或对方材料。该项锚定的是后续追踪问题，影响范围须按地区与群体分别判断。来源保留：战后、改革后或条约后的记忆、地方记录和官方总结常具有事后选择性；后续影响须按地点、群体和时间分层，不作整体化推断。
\eventanchor{EQ-1912-QING-ABDICATION-PRELUDE}{「宣统帝颁布退位诏书，清帝国…}宣统四年（1912年），「宣统帝颁布退位诏书，清帝国终结」之前的条件、信息和争议已通过中央与地方材料显现，构成该事件可追踪的前史节点。核心取证：《宣统政纪》宣统四年相关条；宫中朱批奏折、内阁大库题本与《清史稿》相关志传；同时期地方档案、地方志、日记或对方材料应按具体地区补检。该项锚定的是前史条件与信息背景，不把条件直接写成唯一因果。来源保留：官修编年和地方材料的记录密度与立场并不相同；本行不将条件、传闻、呈报或后见解释直接等同为确定因果。
\eventanchor{EQ-1912-QING-ABDICATION-DECISION}{围绕「宣统帝颁布退位诏书，清…}宣统四年（1912年），围绕「宣统帝颁布退位诏书，清帝国终结」的命令、调度、照会或呈报形成独立的中央—地方行动文书链。核心取证：《宣统政纪》宣统四年相关条；宫中朱批奏折、内阁大库题本与《清史稿》相关志传。该项锚定的是命令或决策环节，命令抵达和结果仍须另外核验。来源保留：奏折、上谕、部议、军机档或条约可证明行政动作和机构立场，不自动证明所有地区、群体或对手按其意图行动。
\eventanchor{EQ-1912-QING-ABDICATION-AFTERMATH}{「宣统帝颁布退位诏书，清帝国…}宣统四年（1912年），「宣统帝颁布退位诏书，清帝国终结」引发的善后、执行或地方回应构成可由不同材料追踪的后续节点。核心取证：《宣统政纪》宣统四年相关条；宫中朱批奏折、内阁大库题本与《清史稿》相关志传；相关地区的档案、志书、契约、报刊或对方材料。该项锚定的是后续追踪问题，影响范围须按地区与群体分别判断。来源保留：战后、改革后或条约后的记忆、地方记录和官方总结常具有事后选择性；后续影响须按地点、群体和时间分层，不作整体化推断。
